% ====================================================================
%  GUIA DE ESTUDIO COMPLETA -- Proyecto #7 (MLP paralelo)
%  Computacion Paralela y Distribuida (CS4052)
%  Lee este documento de principio a fin. No asume ningun conocimiento previo.
%  Compila con:  tectonic guia_completa.tex
% ====================================================================
\documentclass[11pt,a4paper]{article}
\usepackage[utf8]{inputenc}
\usepackage[T1]{fontenc}
\usepackage[spanish,es-tabla]{babel}
\usepackage[a4paper,margin=2.3cm]{geometry}
\usepackage{amsmath,amssymb}
\usepackage{xcolor}
\usepackage{booktabs}
\usepackage{hyperref}
\usepackage{tcolorbox}
\tcbuselibrary{breakable}
\hypersetup{colorlinks=true,linkcolor=blue!55!black,urlcolor=blue!55!black}
\setlength{\parindent}{0pt}\setlength{\parskip}{5pt}

\definecolor{clave}{RGB}{0,90,160}
\definecolor{fondo}{RGB}{245,248,252}
\newtcolorbox{caja}[1][]{
  colback=fondo,colframe=clave!60,boxrule=0.5pt,arc=2pt,
  left=6pt,right=6pt,top=4pt,bottom=4pt,breakable,#1}
\newtcolorbox{importante}[1][]{
  colback=yellow!8,colframe=orange!70!black,boxrule=0.6pt,arc=2pt,
  left=6pt,right=6pt,top=4pt,bottom=4pt,breakable,#1}
\newcommand{\term}[1]{\textbf{\textcolor{clave}{#1}}}

\title{\Huge\textbf{Guía de Estudio Completa}\\[6pt]
\Large Proyecto \#7: Redes Neuronales --- Entrenamiento paralelo de MLP\\[4pt]
\large Computación Paralela y Distribuida (CS4052) -- 2026-I\\[2pt]
\small \emph{Todo explicado desde cero, sin asumir conocimientos previos}}
\author{Kalos Lazo Mera \quad Aaron C. A. Navarro Mendoza \quad Fernando A. Usurin Arias}
\date{Julio de 2026}

\begin{document}
\maketitle
\tableofcontents
\newpage

% ====================================================================
\begin{caja}[title=\textbf{Cómo usar esta guía}]
Este documento explica \textbf{absolutamente todo} lo que necesitas saber para
entender el proyecto y defender la exposición oral. Está ordenado de lo más
básico a lo específico. Si lees de principio a fin, entenderás:
\begin{enumerate}
  \item Qué es la computación paralela y sus conceptos (Parte I).
  \item El modelo teórico PRAM que usa el curso (Parte II).
  \item Qué es una red neuronal / MLP (Parte III).
  \item El problema concreto del proyecto (Parte IV).
  \item El algoritmo paralelo que diseñamos (Parte V).
  \item Cómo se calculan las métricas, con números reales (Parte VI).
  \item Cómo está hecho el código (Parte VII).
  \item Qué significan los resultados (Parte VIII).
  \item Cómo exponer y qué preguntará el profe (Parte IX).
\end{enumerate}
Las cajas \textcolor{orange!70!black}{naranjas} marcan lo MÁS IMPORTANTE.
\end{caja}

% ====================================================================
\part*{PARTE I --- Fundamentos de computación paralela}
\addcontentsline{toc}{section}{PARTE I --- Fundamentos de computación paralela}

\section{¿Qué es la computación paralela?}

Un programa \term{secuencial} ejecuta una instrucción tras otra, en un solo
procesador, una a la vez. Si una tarea tarda 10 horas, así se queda.

La \term{computación paralela} consiste en usar \textbf{varios procesadores a la
vez} para resolver un problema \textbf{más rápido}. La idea es dividir el trabajo
en partes independientes y que cada procesador haga una parte al mismo tiempo
(\emph{en paralelo}).

\textbf{Ejemplo cotidiano.} Si tienes que lavar 100 platos y lo haces solo,
tardas 100 minutos. Si invitas a 4 amigos y cada uno lava 25 plados, terminan en
$\approx 25$ minutos. Eso es paralelismo: el trabajo total es el mismo, pero lo
repartiste.

\begin{importante}[title=\textbf{Idea clave}]
Paralelizar NO reduce el trabajo total (los 100 plados se lavan igual). Reduce
el \textbf{tiempo de reloj} (cuánto tardas tú en ver el resultado) porque varias
unidades de trabajo ocurren \emph{al mismo tiempo}.
\end{importante}

\section{Procesador, núcleo, hilo, proceso}

\begin{itemize}
  \item \term{CPU / procesador}: el ``cerebro'' del computador que ejecuta
        instrucciones.
  \item \term{Núcleo (core)}: una unidad de cálculo dentro de la CPU. Una CPU
        moderna tiene \textbf{varios núcleos} (tu MacBook M5 tiene 10; un nodo del
        cluster Khipu tiene 32). Cada
        núcleo puede ejecutar una secuencia de instrucciones a la vez.
  \item \term{Hilo (thread)}: una secuencia de instrucciones en ejecución. Un
        núcleo puede cambiar rápido entre hilos, pero realmente hace uno a la vez
        (o dos con \emph{hyper-threading}).
  \item \term{Proceso}: un programa en ejecución, con su propia memoria. Un
        proceso puede tener varios hilos dentro.
\end{itemize}

En este proyecto usamos \textbf{procesos} (no hilos) para paralelizar. Verás por
qué en la Parte V (el GIL de Python).

\section{Memoria compartida vs memoria distribuida}

Esta es LA distinción más importante del curso. Hay dos formas de organizar una
máquina paralela:

\subsection{Memoria compartida (shared memory)}
Todos los procesadores/núcleos ven \textbf{la misma memoria}. Si un procesador
escribe un dato en la dirección 1000, los demás lo pueden leer de ahí
inmediatamente. Comunicarse entre procesadores = leer/escribir variables
compartidas. Rápido y simple.

\textbf{Tecnologías:} \term{OpenMP} (en C/C++/Fortran), \term{multithreading},
\term{multiprocessing} (Python). Tu CPU multinúcleo es de memoria compartida.

\subsection{Memoria distribuida (distributed memory)}
Cada procesador tiene \textbf{su propia memoria privada}. Para que el procesador
B vea un dato del procesador A, A debe \emph{enviarle un mensaje} (copiar los
datos por la red). Comunicarse = enviar y recibir mensajes.

\textbf{Tecnología:} \term{MPI} (Message Passing Interface). Los clústers de
muchas computadoras son de memoria distribuida.

\begin{caja}[title=\textbf{Comparación rápida}]
\begin{tabular}{p{0.32\textwidth}p{0.31\textwidth}p{0.31\textwidth}}
& \textbf{Compartida} & \textbf{Distribuida} \\\midrule
Memoria & Una sola, global & Una por procesador \\
Comunicación & Leer/escribir variables & Enviar mensajes (MPI) \\
Ejemplo & Tu laptop multinúcleo & Clúster de PCs \\
Librería típica & OpenMP & MPI \\
¿Copiar datos? & No, ya están & Sí, hay que enviarlos \\
\end{tabular}
\end{caja}

\textbf{En nuestro proyecto usamos memoria compartida.} La justificación está en
la Parte V (los datos son compartidos y de solo lectura, así que es más
eficiente no tener que copiarlos).

\section{Tiempo secuencial y tiempo paralelo}

\begin{itemize}
  \item \term{$T_s$} (tiempo secuencial): cuánto tarda el programa con
        \textbf{1} solo procesador. Es tu línea base.
  \item \term{$T_p$} (tiempo paralelo): cuánto tarda con \textbf{$p$}
        procesadores.
\end{itemize}

El objetivo siempre es que $T_p < T_s$, y cuanto menor, mejor.

\section{Speedup (aceleración)}

El \term{speedup} mide cuántas veces más rápido es el paralelo:

\[
S(p)=\frac{T_s}{T_p}
\]

\textbf{Ejemplo con datos reales de nuestro proyecto} (Khipu, n=40000 muestras):
\begin{itemize}
  \item $T_s = T_1 = 305$\,s (con 1 proceso).
  \item Con $p=8$ procesos: $T_8 = 53$\,s.
  \item $S(8)=305/53 = 5.7$. O sea, 5.7 veces más rápido.
\end{itemize}

El \term{speedup ideal} es $S(p)=p$ (el doble de procesadores $\Rightarrow$ el
doble de rápido). En la práctica casi nunca se alcanza.

\section{Eficiencia}

La \term{eficiencia} es el speedup ``por procesador'':

\[
E(p)=\frac{S(p)}{p}=\frac{T_s}{p\cdot T_p}
\]

Mide qué \emph{tan bien} estás aprovechando los procesadores:
\begin{itemize}
  \item $E=1$ (100\,\%): aprovechamiento perfecto.
  \item $E=0.5$: estás usando la mitad del potencial.
\end{itemize}

\textbf{Ejemplo real:} con $p=4$, $S(4)=3.5$, entonces $E(4)=3.5/4=0.87$. Estás
aprovechando el 87\,\% de los 4 procesadores. ¡Bastante bueno!

\section{La ley de Amdahl}

\begin{importante}[title=\textbf{Ley de Amdahl --- la idea central}]
Ningún programa es 100\,\% paralelizable. Siempre hay una \textbf{fracción
secuencial} $f$ que no se puede paralelizar (leer datos, inicializar, reducir el
resultado final\ldots). Esa fracción fija un \textbf{techo} al speedup por muchos
procesadores que agregues.
\end{importante}

Si una fracción $f$ del programa es secuencial y el resto $(1-f)$ es paralela,
entonces:

\[
S(p) \le \frac{1}{f+\dfrac{1-f}{p}}
\]

Cuando $p\to\infty$ (procesadores infinitos), $S\to 1/f$. \textbf{El speedup
máximo es $1/f$, por muchos núcleos que tengas.}

\textbf{Ejemplo:} si $f=0.1$ (10\,\% secuencial), el speedup máximo es
$1/0.1=10$, aunque tengas 1000 núcleos. Esa fracción secuencial te condena.

En nuestro proyecto, $S$ satura en \textbf{$\approx 5$ en la Mac (10 núcleos) y
$\approx 10.6$ en Khipu (32 núcleos)}. Para la Mac, eso implica una fracción
efectiva $f\approx 1/5 = 0.2$. ¿De dónde viene ese 20\,\% ``secuencial''? No es
código secuencial puro, sino que \textbf{los procesos se estorban} (contención de
memoria) y el hyper-threading infla el conteo de núcleos. Lo veremos en la
Parte VIII.

\section{La ley de Gustafson (weak scaling)}

Amdahl asume que el \textbf{tamaño del problema es fijo}. Gustafson dice: ¿y si
cuando agregas procesadores también \textbf{agrandas el problema}? (cada
procesador hace el mismo trabajo pero en total se resuelve algo más grande).

Bajo este ``weak scaling'', el speedup puede crecer \emph{casi lineal} con $p$,
porque la fracción secuencial deja de dominar.

\begin{caja}[title=\textbf{Resumen Amdahl vs Gustafson}]
\begin{tabular}{lll}
& \textbf{Amdahl} & \textbf{Gustafson} \\\midrule
Tamaño del problema & Fijo & Crece con $p$ \\
Tipo de escalabilidad & Strong scaling & Weak scaling \\
Speedup & Techo limitado por $f$ & Casi lineal \\
Eficiencia $E$ & Cae con $p$ & Se mantiene $\approx 1$ \\
\end{tabular}
\end{caja}

\section{Strong scaling vs weak scaling}

\begin{itemize}
  \item \term{Strong scaling}: mantienes el problema \textbf{fijo} (mismos datos)
        y aumentas $p$. Pregunta: ``¿termina más rápido?'' Mide $T_p$ vs $p$.
        Ideal: $T_p = T_s/p$.
  \item \term{Weak scaling}: aumentas el problema \textbf{junto con} $p$ (cada
        procesador siempre hace el mismo trabajo). Pregunta: ``¿sigue tardando lo
        mismo?'' Mide $T_p(n\cdot p)/T_p(n)$; ideal $=1$.
\end{itemize}

\section{FLOPs y GFLOP/s}

\begin{itemize}
  \item \term{FLOP} = \emph{Floating-point OPeration}: una operación de coma
        flotante (una suma o multiplicación con decimales).
  \item \term{FLOPs} (con minúscula): cantidad total de operaciones.
  \item \term{FLOP/s} o \term{GFLOP/s}: operaciones por segundo (giga =
        $10^9$). Mide el \textbf{rendimiento}: cuán rápido calcula la máquina.
\end{itemize}

\textbf{Fórmula:} $\text{GFLOP/s} = \dfrac{\text{FLOPs totales}}{T_p \cdot 10^9}$.

A más GFLOP/s, más rápido calcula. En nuestro proyecto, el rendimiento pico es
$\approx 53$\,GFLOP/s.

\begin{importante}[title=\textbf{Las 3 métricas que pide el enunciado}]
\begin{enumerate}
  \item \textbf{Tiempo de ejecución} $T_p$ (segundos).
  \item \textbf{Speedup} $S(p)=T_s/T_p$.
  \item \textbf{Eficiencia} $E(p)=S(p)/p$.
\end{enumerate}
Y de regalo: GFLOP/s (rendimiento).
\end{importante}

% ====================================================================
\part*{PARTE II --- El modelo PRAM}
\addcontentsline{toc}{section}{PARTE II --- El modelo PRAM}

\section{¿Qué es PRAM?}

\term{PRAM} = \emph{Parallel Random Access Machine}. Es un \textbf{modelo
teórico} (como una receta abstracta) para describir algoritmos paralelos. No es
una máquina real; es una forma de \textbf{escribir y analizar} algoritmos
paralelos de manera limpia.

Imagina una máquina con:
\begin{itemize}
  \item Un número $p$ de \textbf{procesadores} idénticos: $P_0, P_1, \ldots,
        P_{p-1}$.
  \item Una \textbf{memoria global compartida} con celdas.
  \item Un \textbf{reloj común}: en cada ``paso'' de reloj, todos los
        procesadores ejecutan una instrucción a la vez.
\end{itemize}

El PRAM es la forma estándar del curso de \textbf{describir un algoritmo
paralelo en pseudocódigo} y calcular su complejidad.

\section{Los 4 sabores de PRAM (EREW, CREW, ERCW, CRCW)}

El problema clave de la memoria compartida: \textbf{¿qué pasa si dos
procesadores quieren acceder a la MISMA celda de memoria al mismo tiempo?}

Depende de si es lectura o escritura:
\begin{itemize}
  \item \textbf{Lectura concurrente}: ¿pueden varios leer la misma celda a la vez?
  \item \textbf{Escritura concurrente}: ¿pueden varios escribir en la misma celda a la vez?
\end{itemize}

De ahí salen 4 modelos (E=Exclusive, C=Concurrent):

\begin{caja}[title=\textbf{Los 4 modelos PRAM}]
\begin{tabular}{lll}
Modelo & Lectura & Escritura \\\midrule
\textbf{EREW} & Exclusiva & Exclusiva \\
\textbf{CREW} & Concurrente & Exclusiva \\
\textbf{ERCW} & Exclusiva & Concurrente \\
\textbf{CRCW} & Concurrente & Concurrente \\
\end{tabular}
\end{caja}

\begin{itemize}
  \item \textbf{EREW} (el más estricto): ni leer ni escribir pueden coincidir. El
        más ``caro'' de simular pero el más realista en hardware básico.
  \item \textbf{CREW}: varios pueden LEER a la vez, pero si alguien ESCRIBE, debe
        ser el único. \textbf{Es el modelo que usamos en el proyecto.}
  \item \textbf{CRCW} (el más permisivo): hasta pueden escribir varios a la vez
        (hay que definir cómo se resuelve el conflicto). El más poderoso pero el
        menos realista.
\end{itemize}

\textbf{¿Por qué elegimos CREW en el proyecto?} Porque:
\begin{itemize}
  \item Los $p$ procesadores \textbf{leen el dataset concurrentemente} (lectura
        concurrente $\Rightarrow$ la C de CREW).
  \item Las escrituras del resultado son \textbf{exclusivas}: cada procesador
        escribe en su propia celda, sin pisarse (escritura exclusiva $\Rightarrow$
        la EW de CREW).
  \item No necesitamos escritura concurrente (CRCW sería matar moscas a cañonazos).
\end{itemize}

\section{El paradigma Work--Time (WT)}

El curso usa una forma específica de escribir algoritmos PRAM llamada
\term{Work--Time (WT)}. En vez de contar ciclos de reloj uno por uno, se escribe
el pseudocódigo marcando qué partes van en paralelo, y luego se derivan dos
cantidades:

\begin{itemize}
  \item \term{Work $W$}: el trabajo total = número de operaciones = equivale al
        tiempo secuencial $T_s$.
  \item \term{Span $T_\infty$}: el camino crítico = el tiempo con procesadores
        \emph{infinitos}.
\end{itemize}

\section{La palabra clave \texttt{pardo}}

En el pseudocódigo WT, la palabra \term{\texttt{pardo}} (\emph{parallel-do})
marca los bucles que se ejecutan \textbf{en paralelo}. Todo lo que está dentro de
un \texttt{pardo} lo hacen todos los procesadores a la vez.

\begin{verbatim}
for k = 0 to p-1 pardo do
    <lo que hace cada procesador, en paralelo>
\end{verbatim}

\begin{importante}[title=\textbf{Corrección clave del parcial}]
En el informe parcial llamamos \textbf{``Broadcast''} al paso donde los
procesadores leen los datos. \textbf{Eso estaba mal}: en memoria compartida NO
existe un ``broadcast''. Los datos ya están en memoria global; cada procesador
simplemente los \textbf{lee concurrentemente} (eso es justamente el acceso CREW).
Por eso en la versión final lo renombramos a \textbf{``lectura concurrente''} y
lo marcamos con \texttt{pardo} donde corresponde. \textbf{No digas nunca
``broadcast'' en memoria compartida} o el profe te para.
\end{importante}

\section{Work $W$ y Span $T_\infty$ con detalle}

\subsection{Work $W$}
Es el \textbf{total de operaciones} que hace el algoritmo, contando las de todos
los procesadores. Si lo corriera un solo procesador secuencialmente, tardaría
$O(W)$. Por eso \textbf{$W$ equivale al tiempo secuencial $T_s$}.

\subsection{Span $T_\infty$}
Es la duración del \textbf{camino crítico}: la cadena más larga de operaciones
que dependen una de otra. Representa el tiempo que tardarías si tuvieras
\textbf{procesadores infinitos} (por eso el subíndice $\infty$). Con procesadores
infinitos, lo único que limita es qué operaciones DEBEN ir en secuencia.

\textbf{Intuición:} imagina una receta. Aunque tuvieras 1000 cocineros, no
puedes servir el plato antes de hornearlo. El span es el tiempo mínimo
irreducible (lo que forzosamente va en secuencia).

\section{El Teorema de Brent}

El Teorema de Brent relaciona $W$, $T_\infty$ y el número real de procesadores
$p$ con el tiempo paralelo $T_p$:

\[
\boxed{T_p \;\le\; \frac{W}{p} + T_\infty}
\]

\textbf{Interpretación:}
\begin{itemize}
  \item $\dfrac{W}{p}$: si repartieras el trabajo $W$ perfectamente entre $p$
        procesadores.
  \item $T_\infty$: el ``mínimo irreducible'' (camino crítico).
  \item El tiempo real está por debajo de la suma.
\end{itemize}

Esto da una cota \textbf{superior} del tiempo paralelo: en el peor caso, $T_p$
no supera $W/p + T_\infty$.

\section{Costo $C_p$ y WT-optimalidad}

\term{Costo} $C_p = T_p \cdot p$ (tiempo $\times$ número de procesadores). Es una
medida del ``esfuerzo total'' gastado.

Un algoritmo PRAM es \term{WT-óptimo} si su costo es proporcional al trabajo:
\[
C_p = O(W) \quad\Longleftrightarrow\quad T_p = O(W/p)
\]
Es decir, cuando el paralelo reparte \emph{perfectamente} el trabajo (no desperdicia
procesadores). Esto se cumple si $p \le W/T_\infty$.

\begin{caja}[title=\textbf{En nuestro proyecto}]
$W/T_\infty = O(p_s)$ (número de semillas). Como siempre usamos $p \le p_s=32$,
el algoritmo \textbf{es WT-óptimo}. Bien.
\end{caja}

% ====================================================================
\part*{PARTE III --- Redes neuronales (MLP) desde cero}
\addcontentsline{toc}{section}{PARTE III --- Redes neuronales (MLP) desde cero}

\section{¿Qué es una red neuronal?}

Una \term{red neuronal artificial} es un modelo matemático inspirado (muy
vagamente) en las neuronas del cerebro. Sirve para \textbf{aprender patrones a
partir de datos}: le muestras muchos ejemplos y ``aprende'' a hacer una tarea
(clasificar imágenes, predecir un número, etc.).

No es magia ni inteligencia: es un montón de \textbf{multiplicaciones y sumas}
organizadas de forma ingeniosa, que se ajustan automáticamente para minimizar
errores.

\section{La neurona artificial}

Una \term{neurona} (o ``perceptrón'') es una función matemática simple:
\[
y = f(w_1 x_1 + w_2 x_2 + \ldots + w_d x_d + b)
\]
\begin{itemize}
  \item $x_1,\ldots,x_d$: las \textbf{entradas} (los datos, ej. los píxeles de
        una imagen o las ``features'').
  \item $w_1,\ldots,w_d$: los \textbf{pesos} (números que la red aprende).
  \item $b$: el \textbf{sesgo} (bias), otro número que se aprende.
  \item $f$: la \textbf{función de activación} (veremos cuál).
  \item $y$: la \textbf{salida}.
\end{itemize}

La neurona calcula una \textbf{suma ponderada} de las entradas (cada entrada por
su peso) más el sesgo, y le aplica una función.

\section{MLP: la estructura}

\term{MLP} = \emph{Multi-Layer Perceptron} = perceptrón multicapa. Son varias
neuronas organizadas en \textbf{capas}:

\begin{verbatim}
[ entrada ]  ->  [ capa oculta ]  ->  [ salida ]
  d neuronas     h neuronas          c neuronas
\end{verbatim}

\begin{itemize}
  \item \textbf{Capa de entrada}: recibe los datos. Tiene $d$ valores (en nuestro
        proyecto $d=32$ ``features'').
  \item \textbf{Capa oculta}: las neuronas ``intermedias'' que hacen el trabajo
        pesado. Tiene $h$ neuronas (en nuestro proyecto $h=128$).
  \item \textbf{Capa de salida}: produce el resultado. Tiene $c$ neuronas (en
        nuestro proyecto $c=10$ clases).
\end{itemize}

Cada neurona de una capa está \textbf{conectada} a todas las de la siguiente, y
cada conexión tiene un peso. De ahí sale el costo computacional (multiplicar
matrices).

\section{Forward propagation (propagación hacia adelante)}

Es el cálculo de la salida a partir de la entrada. Paso a paso:

\textbf{De la entrada a la capa oculta:} cada neurona oculta $j$ calcula:
\[
z_j = \sum_{i=1}^{d} w_{ji}^{(1)}\, x_i + b_j^{(1)}, \qquad a_j = f(z_j)
\]
Son $h$ neuronas, cada una hace $d$ multiplicaciones. Total: $h\cdot d$
multiplicaciones. Con activación.

\textbf{De la capa oculta a la salida:} cada neurona de salida $k$ calcula:
\[
o_k = \sum_{j=1}^{h} w_{kj}^{(2)}\, a_j + b_k^{(2)}
\]
Son $c$ neuronas, cada una hace $h$ multiplicaciones. Total: $c\cdot h$.

\textbf{Total de operaciones por una muestra:} $\approx 2\cdot h\cdot d + 2\cdot
c\cdot h$ (contando multiplicaciones y sumas como FLOPs).

\section{Función de activación}

La función $f$ introduce \textbf{no linealidad} (sin ella, una red de varias
capas sería equivalente a una sola, y no podría aprender cosas complejas). Las
típicas:
\begin{itemize}
  \item \term{ReLU}: $f(z)=\max(0,z)$. Simple y muy usada.
  \item \term{Sigmoide}: aplasta a $[0,1]$.
  \item \term{Softmax} (en la salida de clasificación): convierte las salidas en
        probabilidades (suman 1).
\end{itemize}

\section{Función de pérdida (loss)}

Después del forward, la red produce una predicción. La comparamos con la
respuesta correcta y medimos el \textbf{error} con una función de pérdida. Para
clasificación se usa típicamente la \term{entropía cruzada} (\emph{cross-entropy}):
\[
\text{Loss} = -\sum_k y_k \log(\hat{y}_k)
\]
donde $y$ es la respuesta correcta y $\hat{y}$ la predicción. Menor pérdida =
mejor predicción.

\section{Backpropagation (retropropagación)}

Aquí está la magia del aprendizaje. La red quiere \textbf{reducir la pérdida}.
Para eso necesita saber \textbf{cómo cambiar cada peso} para que el error baje.

\term{Backpropagation} calcula, usando la regla de la cadena del cálculo, el
\textbf{gradiente} de la pérdida respecto a cada peso: $\partial \text{Loss}/
\partial w$. Es decir, ``cuánto contribuye cada peso al error''.

\textbf{Costo:} el backward cuesta más o menos lo mismo que el forward (hay que
recorrer la red al revés). Por eso decimos que entrenar cuesta $\approx 2\times$
un forward.

\section{Descenso de gradiente}

Una vez que sabes el gradiente, \textbf{ajustas los pesos} en dirección contraria
al error:
\[
w_{\text{nuevo}} = w_{\text{viejo}} - \eta \cdot \frac{\partial \text{Loss}}{\partial w}
\]
\begin{itemize}
  \item $\eta$ (\term{learning rate}, tasa de aprendizaje): qué tan grande es el
        paso. Si es muy grande, diverge; muy pequeño, aprende lento.
\end{itemize}
Repetido muchas veces, los pesos convergen a valores que minimizan el error.

\section{Épocas, learning rate, batch}

\begin{itemize}
  \item \term{Época (epoch)}: una pasada completa por \emph{todos} los datos de
        entrenamiento. En nuestro proyecto $E=100$ épocas.
  \item \term{Batch}: en la práctica se procesan los datos en grupos (lotes) en
        vez de uno por uno. Nosotros usamos todo el conjunto.
  \item \term{Iteraciones}: el número de actualizaciones. Una época con $n$
        muestras y batch $b$ tiene $n/b$ iteraciones.
\end{itemize}

\section{Accuracy (exactitud)}

Métrica de calidad: fracción de predicciones correctas. Si de 100 ejemplos
acierta 80, $accuracy=0.80$. En nuestro proyecto los mejores modelos alcanzan
$\approx 0.77$--$0.87$ según el tamaño de los datos.

\section{Los hiperparámetros del proyecto}

Un \term{hiperparámetro} es una configuración que tú eliges (no la aprende la
red). Los del proyecto:
\begin{itemize}
  \item $h=128$: neuronas de la capa oculta.
  \item $E=100$ (\texttt{max\_iter}): épocas máximas.
  \item $\alpha=10^{-4}$: regularización L2 (evita sobreajuste).
  \item $\eta=10^{-3}$ (\texttt{learning\_rate\_init}): learning rate.
  \item $d=32$: features de entrada; $c=10$: clases.
\end{itemize}

\section{Las semillas y por qué entrenar varios MLP}

\begin{importante}[title=\textbf{La pregunta clave del proyecto}]
El entrenamiento de una red empieza con pesos \textbf{aleatorios}, generados por
un \term{generador de números pseudoaleatorios} (RNG) que se controla con una
\textbf{semilla} (\emph{seed}). \textbf{Cambiar la semilla cambia el resultado}:
distintos pesos iniciales $\Rightarrow$ distinto modelo $\Rightarrow$ distinta
exactitud.

Como no sabemos de antemano qué semilla da el mejor modelo, una práctica robusta
es \textbf{entrenar varios MLP con semillas distintas y conservar el de mayor
exactitud}. \textbf{¡Eso es justo lo que paralelizamos!}
\end{importante}

Entrenamos $p_s=32$ MLP, cada uno con una semilla $s\in\{0,1,\ldots,31\}$, y nos
quedamos con el de mejor accuracy. Y como cada MLP es \textbf{independiente} de
los demás, los podemos entrenar \textbf{en paralelo}.

% ====================================================================
\part*{PARTE IV --- El proyecto}
\addcontentsline{toc}{section}{PARTE IV --- El proyecto}

\section{El enunciado (proyecto \#7)}

El proyecto pide: dado el pseudocódigo secuencial del Algoritmo 1 (entrenar
varios MLP con distintas semillas y quedarse con el mejor), \textbf{paralelizarlo}.

Requisitos del enunciado:
\begin{enumerate}
  \item Hacer un \textbf{PRAM} y el código en paralelo.
  \item Usar $p\in\{1,2,4,8,16,32\}$ procesos y $n\in\{5000,10000,20000,40000\}$
        muestras; generar gráficas de tiempos.
  \item Registrar el desarrollo en \textbf{al menos 3 versiones beta}.
  \item Comparar los resultados con la \textbf{complejidad teórica}, superponiendo
        la curva teórica a la experimental.
  \item Calcular \textbf{speedup y eficiencia}, analizar la escalabilidad y
        responder: ¿se mantiene $E$ constante?
\end{enumerate}

\section{El algoritmo secuencial (Algoritmo 1)}

\begin{caja}[title=\textbf{Algoritmo 1 --- secuencial}]
\begin{verbatim}
Dividir (X, y) en train y test  (una sola vez)
mejor_accuracy = -inf
para cada semilla s en {0,...,31}:
    crear modelo M con semilla s
    entrenar M con (X_train, y_train)     <- esto es lo pesado
    A = accuracy(M, X_test, y_test)
    si A > mejor_accuracy:
        mejor_accuracy = A
        mejor_modelo = M
devolver mejor_modelo
\end{verbatim}
\end{caja}

El bucle entrena 32 modelos \textbf{uno detrás de otro}. Lento. Ese es nuestro
$T_s$.

\section{¿Por qué es ``embarassingly parallel''?}

\begin{importante}[title=\textbf{Embarassingly parallel = vergonzosamente paralelo}]
Es un término técnico (en serio) para problemas que se paralelizan de forma
\textbf{trivial}: el trabajo se divide en partes \textbf{totalmente
independientes}, sin necesidad de comunicación entre ellas durante el cómputo.

El entrenamiento de los 32 MLP es así: \textbf{cada MLP no necesita nada de los
demás}. Solo se juntan al final para comparar accuracies y elegir el mejor
(una reducción \emph{arg-máx}). Por eso paraleliza genial.
\end{importante}

% ====================================================================
\part*{PARTE V --- El algoritmo paralelo}
\addcontentsline{toc}{section}{PARTE V --- El algoritmo paralelo}

\section{La idea del algoritmo paralelo}

En vez de entrenar los 32 MLP uno por uno, los \textbf{repartimos} entre $p$
procesadores. Cada procesador entrena $\lceil 32/p \rceil$ MLP y reporta su
mejor modelo local. Al final, los $p$ ganadores locales se comparan para sacar
el global (reducción \emph{arg-máx} en árbol).

\section{El modelo PRAM CREW (corregido)}

\begin{caja}[title=\textbf{Algoritmo 2 --- PRAM CREW (memoria compartida)}]
\begin{verbatim}
Paso 1 (secuencial, raíz):
    Dividir (X,y) en train/test
    Dejar los datos en MEMORIA COMPARTIDA

Paso 2 (PARALELO, pardo):   <- lo que antes llamábamos "Broadcast"
    para k=0..p-1 pardo:
        P_k LEE de memoria compartida (X_train, ...)   [lectura CREW]
        P_k entrena sus semillas S_k y guarda su mejor local

Pasos 3..(2+log2 p) (PARALELO, pardo):  reducción arg-máx en árbol
    para stride=1..log2(p):
        para k pardo:
            si mejor[k+stride] > mejor[k]: mejor[k] = mejor[k+stride]

Devolución: mejor[0] = M*
\end{verbatim}
\end{caja}

\section{Las correcciones que hicimos al parcial}

\begin{enumerate}
  \item \textbf{Quitar ``Broadcast''.} En memoria compartida los datos ya están
        en memoria global; los procesadores los \textbf{leen concurrentemente}
        (acceso CREW). No hay nada que ``difundir''.
  \item \textbf{No paralelizar la lectura como paso aparte.} Es solo lectura de
        referencias ya compartidas.
  \item \textbf{No medir la lectura ni el split en los tiempos.} Medimos solo la
        región paralelizable (entrenamiento) + la reducción.
  \item \textbf{Justificar el uso de memoria compartida} (sección siguiente).
  \item \textbf{Incluir el término $\log_2 p$} en la fórmula de $T_p$ (viene de
        la reducción en árbol).
\end{enumerate}

\section{Justificación de la memoria compartida (¡pregunta favorita del profe!)}

\begin{importante}[title=\textbf{¿Por qué memoria compartida y no MPI?}]
\begin{enumerate}
  \item \textbf{La comunicación es mínima.} Cada MLP es independiente; la única
        interacción es la reducción final (un árbol de $\log_2 p$ rondas). No se
        aprovecha nada del paso de mensajes.
  \item \textbf{Los datos son compartidos y de solo lectura.} Cada MLP se entrena
        con \emph{todo} el dataset. En memoria compartida, los $p$ procesadores
        simplemente \textbf{leen} las mismas referencias (sin copiar). En MPI
        habría que \textbf{copiar el dataset a cada proceso} (memoria $\times p$ +
        tiempo de envío).
  \item \textbf{Conclusión:} la memoria compartida es aquí \textbf{más
        eficiente}. MPI no aportaría nada y sumaría overhead.
\end{enumerate}
\end{importante}

\section{OpenMP, multiprocessing y el GIL}

\begin{itemize}
  \item \term{OpenMP}: la forma estándar de paralelismo en memoria compartida en
        C/C++/Fortran. Usas directivas (\texttt{\#pragma omp parallel}) para
        decir ``este bucle va en paralelo''. Modelo \emph{fork-join}.
  \item \term{GIL} (Global Interpreter Lock): un candado del intérprete de Python
        que impide que \textbf{dos hilos} ejecuten código Python a la vez. Por
        eso en Python los \textbf{hilos} no dan paralelismo real.
  \item \term{multiprocessing}: la solución de Python: en vez de hilos, usa
        \textbf{procesos} separados (cada uno con su intérprete y su GIL). Así sí
        hay paralelismo real. Es el análogo de OpenMP en Python.
\end{itemize}

\textbf{En el proyecto:} el MODELO conceptual es OpenMP (memoria compartida,
fork-join), pero la IMPLEMENTACIÓN es \texttt{multiprocessing} porque el GIL
impide hilos reales. Conceptualmente es lo mismo: $p$ trabajadores que comparten
memoria y se sincronizan al final.

\section{Fork-join}

Es el patrón de ejecución:
\begin{enumerate}
  \item Un proceso ``raíz'' está corriendo.
  \item \textbf{Fork}: el raíz crea $p$ procesos trabajadores (heredan la
        memoria por \emph{copy-on-write}, sin copiar).
  \item Los $p$ trabajadores hacen su parte \textbf{en paralelo}.
  \item \textbf{Join}: el raíz espera a que todos terminen y recoge los
        resultados.
\end{enumerate}
Es exactamente lo que hace \texttt{\#pragma omp parallel} en OpenMP y
\texttt{multiprocessing.Pool} en Python.

\section{fork vs spawn (detalle del código)}

\begin{itemize}
  \item \term{fork}: el proceso hijo es un \textbf{clon} del padre (hereda la
        memoria por copy-on-write, sin copiar datos). \textbf{Eficiente}. Es lo
        fiel al modelo de memoria compartida. Lo usamos en la beta 3.
  \item \term{spawn}: el proceso hijo arranca \textbf{desde cero} y hay que
        \textbf{serializar} (con \emph{pickle}) y enviarle los datos. Más lento.
        La beta 2 lo usaba (y por eso era ineficiente).
\end{itemize}

% ====================================================================
\part*{PARTE VI --- Las métricas, paso a paso (con números reales)}
\addcontentsline{toc}{section}{PARTE VI --- Las métricas, paso a paso}

Aquí derivamos cada métrica del pseudocódigo, con un ejemplo numérico real.

\section{Setup numérico}

Plataforma: \textbf{cluster Khipu, nodo n003} (Intel Xeon Gold 6130, 32 núcleos).
Tomemos $n=40000$ muestras. Parámetros: $p_s=32$ semillas, $E=100$ épocas,
$d=32$ features, $h=128$ neuronas, $c=10$ clases.

\section{Work $W$ (trabajo total)}

El paso dominante es entrenar los 32 MLP. Cada MLP cuesta $O(E\cdot n\cdot
d\cdot h)$:
\[
W = O(p_s \cdot E \cdot n \cdot d \cdot h) = O(32\cdot 100\cdot 40000\cdot 32\cdot 128).
\]
Esto \textbf{coincide con $T_s$}, el tiempo secuencial. Medido en Khipu:
$T_s=307$\,s.

\section{Span $T_\infty$}

Con procesadores infinitos, cada semilla tiene su propio procesador ($p=p_s=32$),
así que cada uno entrena un solo MLP:
\[
T_\infty = O(E\cdot n\cdot d\cdot h) + O(\log p_s) = O(E\cdot n\cdot d\cdot h).
\]
Es el tiempo de \textbf{un} MLP más la reducción. Irreducible.

\section{Tiempo paralelo $T_p$}

Por el Teorema de Brent, $T_p \le W/p + T_\infty$. Como $E,d,h$ son constantes
en nuestros experimentos, se absorben en una constante $a$ [s/muestra]. El modelo
numérico que ajustamos a los datos es:
\[
\boxed{T_p \approx a\cdot\frac{p_s}{p}\cdot n + b\cdot\log_2 p}
\]
\textbf{Ejemplo} ($n=40000$, $p=2$):
\[
T_2 \approx 2.42\!\times\!10^{-4}\cdot\frac{32}{2}\cdot 40000 + 2.18\cdot\log_2 2
= 154.9 + 2.18 = 157\text{ s}.
\]
El medido fue $153.6$\,s. ¡Coincide! (la pequeñísima diferencia es ruido).

\section{Costo $C_p$ y WT-optimalidad}

$C_p = T_p\cdot p$. Como $p\le p_s=32=W/T_\infty$, el algoritmo es WT-óptimo:
$C_p=O(W)$.

\section{Speedup y eficiencia (ejemplo real, Khipu)}

Para $n=40000$ en Khipu (32 núcleos):
\begin{center}
\begin{tabular}{rrrrr}
\toprule
$p$ & $T_p$ (s) & Speedup $S$ & Eficiencia $E$ & GFLOP/s \\\midrule
1 & 305.3 & 1.00 & 1.00 & 7.2 \\
2 & 153.6 & 1.99 & 0.99 & 14.3 \\
4 & 108.9 & 2.80 & 0.70 & 20.3 \\
8 & 53.2 & 5.73 & 0.72 & 41.4 \\
16 & 33.0 & 9.26 & 0.58 & 66.8 \\
32 & 28.9 & 10.58 & 0.33 & 76.3 \\
\bottomrule
\end{tabular}
\end{center}

\textbf{Qué leer aquí:}
\begin{itemize}
  \item $p=2$ es casi ideal (eficiencia 0.99).
  \item El speedup sigue creciendo hasta $p=32$, donde alcanza $\sim 10.6$.
  \item La eficiencia cae a $p$ alto (a $p=32$, 0.33) por la contención de
        memoria y el hyper-threading (32 lógicos $\approx$ 16 físicos).
  \item El GFLOP/s sube hasta $\sim 76$.
\end{itemize}

\section{Cálculo de FLOPs}

Para UN MLP (forward + backward, $E$ épocas, $n$ muestras, arquitectura $d$-$h$-$c$):
\[
\text{FLOPs por MLP} = 4\cdot E\cdot h\cdot(d+c)\cdot n.
\]
\textbf{Ejemplo} ($n=40000$):
$4\cdot 100\cdot 128\cdot(32+10)\cdot 40000 = 8.6\!\times\!10^{10}$ FLOPs $=86$\,GFLOPs por MLP.
Los 32 MLP suman $\approx 2.75\!\times\!10^{12}$ FLOPs. A $T_8=53.2$\,s (Khipu),
$\text{GFLOP/s}=2.75\!\times\!10^{12}/(53.2\cdot 10^9)\approx 52$. $\checkmark$
(En $p=32$, $T=28.9$\,s $\Rightarrow \approx 76$\,GFLOP/s, el pico.)

% ====================================================================
\part*{PARTE VII --- La implementación}
\addcontentsline{toc}{section}{PARTE VII --- La implementación}

\section{Las 3 versiones beta}

El enunciado pide registrar el desarrollo en $\geq 3$ pasos. Nuestras 3 betas:

\begin{enumerate}
  \item \textbf{Beta 1 (\texttt{beta1\_secuencial.py}):} el Algoritmo 1 puro,
        secuencial. Es la \textbf{referencia} (da $T_s$ y verifica correctitud).
  \item \textbf{Beta 2 (\texttt{beta2\_paralelo.py}):} primer intento paralelo
        con \texttt{spawn}. Funciona pero es \textbf{ineficiente} porque
        serializa los datos a cada worker.
  \item \textbf{Beta 3 (\texttt{beta3\_paralelo\_opt.py}):} la versión final,
        con \textbf{fork} (los workers heredan los datos por copy-on-write, sin
        copiar), reparto balanceado de semillas y medición fina de tiempos.
        \textbf{Esta es la que usamos para experimentar.}
\end{enumerate}

\section{Worker crítico (cómo se mide el tiempo paralelo)}

\begin{importante}[title=\textbf{El tiempo paralelo NO es la suma}]
En paralelo, los $p$ trabajadores corren \textbf{a la vez}. El tiempo total de
esa fase es el del \textbf{trabajador más lento} (el último en terminar), NO la
suma de todos. Si sumaras los tiempos estarías midiendo el \emph{trabajo} $W$, no
el \emph{tiempo} $T_p$.

Por eso medimos \texttt{T\_compute = max(tiempos de los workers)}.
\end{importante}

\section{Cómo se descompone el tiempo}

\begin{itemize}
  \item \texttt{T\_compute}: tiempo del worker crítico (la región paralela).
  \item \texttt{T\_reduce}: overhead de recoger resultados + arg-máx
        (\textbf{medido $<1.5\%$}).
  \item \texttt{T\_total = T\_compute + T\_reduce}.
  \item \textbf{NO medimos} la lectura de datos ni el split (corrección del profe).
\end{itemize}

\section{Reparto balanceado de semillas}

Repartimos las 32 semillas en $p$ grupos de tamaño $\lceil 32/p\rceil$ o
$\lfloor 32/p\rfloor$, para que funcione \textbf{aunque $p$ no divida a 32}. En
la práctica usamos $p\in\{1,2,4,8,16,32\}$ que sí dividen a 32, pero el código es
robusto.

% ====================================================================
\part*{PARTE VIII --- Los resultados, interpretados}
\addcontentsline{toc}{section}{PARTE VIII --- Los resultados, interpretados}

\section{El hallazgo más importante: dos plataformas}

\begin{importante}[title=\textbf{Mismo algoritmo, dos hardwares: el speedup escala}]
Corrimos el mismo código en dos plataformas de memoria compartida:
\textbf{MacBook M5 (10 núcleos)} y \textbf{cluster Khipu, nodo n003 (32 núcleos,
Xeon Gold 6130)}. Resultado: el speedup pasa de \textbf{$\sim$5 (Mac)} a
\textbf{$\sim$10.6 (Khipu)} al aumentar los núcleos. Eso \textbf{valida
experimentalmente} que el algoritmo escala con el hardware, como predice la
teoría. Y en ambas plataformas el speedup no llega a $p$ (los núcleos totales):
investiguemos por qué.
\end{importante}

\section{¿Por qué el speedup no llega a los núcleos totales?}

Descartamos causas con el método científico:

\textbf{Hipótesis 1: ¿es el overhead de gestión de procesos (fork)?}
Medimos \texttt{T\_reduce}: es $<1.5\%$. \textbf{Descartado.}

\textbf{Hipótesis 2: ¿es que sklearn paraleliza cada MLP con BLAS y se estorban?}
Medimos un MLP aislado con y sin limitar hilos: idéntico
(\texttt{threadpoolctl} no detecta hilos). Limitamos el BLAS a 1 hilo por proceso
para que la única paralelización sea la nuestra. \textbf{Descartado.}

\textbf{Hipótesis 3 (correcta): contención de memoria + hyper-threading.}
El entrenamiento del MLP es \emph{memory-bound} (accede a $n\cdot d$ datos por
época). A alto $p$, los procesos compiten por el ancho de banda de memoria y cada
MLP se vuelve más lento. Además, en el Xeon los ``32 núcleos'' son en realidad
\textbf{16 físicos con hyper-threading} (32 lógicos), así que el speedup real
máximo está entre 10 y 16. La Mac, sin hyper-threading, satura antes ($\sim$5 con
sus 10 núcleos).

\begin{caja}[title=\textbf{La lección clave}]
El speedup crece con el hardware (Mac $\sim$5 $\to$ Khipu $\sim$10.6), pero
nunca llega a $p$ porque (1) el hyper-threading infla el conteo de núcleos y (2)
la memoria compartida se satura. Un análisis honesto que mide, descarta
hipótesis y concluye con causas verificadas vale más que un speedup ``perfecto''.
\end{caja}

\section{Validación con la teoría (el ajuste, Khipu)}

Ajustamos el modelo $T_p = a\cdot(p_s/p)\cdot n + b\cdot\log_2 p$ por mínimos
cuadrados no-negativos sobre los 48 puntos de Khipu:
\[
a=2.42\!\times\!10^{-4}\text{ s/muestra},\quad b=2.18\text{ s},\quad R^2=0.989.
\]
\textbf{$R^2=0.989$ significa que el modelo explica el 98.9\,\% de la variación
de los datos.} Un ajuste excelente. La teoría (con el término $\log_2 p$ de la
reducción, que corregimos respecto al parcial) reproduce los datos. (En la Mac el
ajuste dio $R^2=0.982$; ambos excelentes.)

\textbf{Detalle de rigor:} el modelo ajusta bien en la región lineal ($p\le 4$);
para $p>8$ el experimental queda por encima de la teoría \emph{justo por la
contención de memoria} que no captura la fórmula lineal.

\section{¿$E$ se mantiene constante? (pregunta del enunciado)}

\textbf{Respuesta matizada:}
\begin{itemize}
  \item \textbf{Strong scaling} (Amdahl): $E$ \textbf{NO} se mantiene; cae con
        $p$ (por overhead + contención).
  \item \textbf{Weak scaling} (Gustafson): $E$ \textbf{SÍ} se mantiene cercana a
        1 (al crecer $n$ con $p$, la carga por procesador es constante).
\end{itemize}
Lo mostramos en dos gráficas separadas.

% ====================================================================
\part*{PARTE IX --- La exposición oral}
\addcontentsline{toc}{section}{PARTE IX --- La exposición oral}

\section{Estructura recomendada (10 minutos)}

\begin{enumerate}
  \item \textbf{(1 min) Problema:} entrenar 32 MLP con semillas distintas,
        quedarse con el mejor. Es embarazosamente paralelo.
  \item \textbf{(2 min) PRAM CREW:} muestra el Algoritmo 2. Señala \texttt{pardo},
        la \textbf{lectura concurrente} (no ``broadcast''), la reducción arg-máx.
        \textbf{Justifica la memoria compartida}.
  \item \textbf{(2 min) Métricas:} $W$, $T_\infty$, $T_p$ (con el $\log_2 p$),
        speedup, eficiencia, FLOPs. Señala la fórmula con el término de
        reducción.
  \item \textbf{(3 min) Resultados:} dos plataformas (Mac $\sim$5, \textbf{Khipu
        $\sim$10.6}), la comparación valida la escalabilidad. Eficiencia
        (responde ``¿$E$ constante?'' $\to$ strong cae, weak se mantiene),
        validación teórica ($R^2=0.989$), GFLOP/s.
  \item \textbf{(1 min) Mejoras propuestas.}
  \item \textbf{(1 min) Conclusiones.}
\end{enumerate}

\section{Preguntas que hará el profe (con respuestas)}

\textbf{P: ¿Por qué el speedup no llega a los núcleos totales?}\\
R: Dos causas: (1) \textbf{hyper-threading} --- en el Xeon, 32 núcleos lógicos
$\approx$ 16 físicos, así que el speedup real máximo está entre 10 y 16 (medimos
$\sim$10.6); (2) \textbf{contención de memoria} --- el MLP es memory-bound y a
alto $p$ los procesos compiten. La reducción arg-máx es $<1.5\%$, así que NO es
la comunicación. La prueba: en la Mac (10 núcleos, sin HT) satura en $\sim$5, y
en Khipu (32) llega a $\sim$10.6 --- el speedup crece con el hardware.

\textbf{P: ¿Por qué corrieron en dos máquinas?}\\
R: Para validar la escalabilidad en hardware distinto. Mismo código, mismo
dataset, dos plataformas de memoria compartida. Mostrar que el speedup se duplica
al pasar de 10 a 32 núcleos valida experimentalmente que el algoritmo escala.

\textbf{P: ¿Por qué no usaron MPI?}\\
R: El problema es embarazosamente paralelo y los datos son compartidos de solo
lectura. En MPI habría que copiar el dataset a cada proceso (memoria $\times p$).
En memoria compartida cada proceso hereda los datos por fork sin copiar. MPI no
aportaría nada y sumaría overhead.

\textbf{P: ¿Es realmente OpenMP si usan Python?}\\
R: El MODELO es OpenMP (memoria compartida, fork-join). La IMPLEMENTACIÓN es
\texttt{multiprocessing} porque el GIL impide hilos reales. Conceptualmente es
lo mismo.

\textbf{P: ¿Cómo verificaron que el paralelo da lo mismo que el secuencial?}\\
R: El código es determinista (semillas fijas). Comparamos las 32 accuracies y el
modelo ganador entre sec y par: \textbf{idénticos}.

\textbf{P: ¿Qué es el ``worker crítico'' que miden?}\\
R: El tiempo de la región paralela = el worker más lento. Medir la suma sería
medir el trabajo $W$, no el tiempo $T_p$.

\textbf{P: ¿Por qué no midieron la lectura de datos?}\\
R: Por la observación del profesor: no incluir I/O en la medición de tiempos.
Medimos solo la región paralelizable + la reducción.

\textbf{P: ¿De dónde sale el $\log_2 p$ en $T_p$?}\\
R: De la reducción arg-máx en árbol binario: con $p$ resultados locales, hace
$\log_2 p$ rondas para encontrar el máximo global.

\textbf{P: ¿El algoritmo es escalable?}\\
R: Sí. Strong scaling: lineal hasta $p=4$. Weak scaling: $E\approx 1$. El límite
es el ancho de banda de memoria, no el algoritmo.

\textbf{P: ¿El algoritmo es WT-óptimo?}\\
R: Sí. $C_p=T_p\cdot p = O(W)$ porque $p\le p_s = W/T_\infty$.

\textbf{P: ¿Cómo calcularon los FLOPs?}\\
R: A partir de la arquitectura: $4\cdot E\cdot h\cdot(d+c)\cdot n$ por MLP. No
usamos hw-counters; contamos las operaciones del costo conocido del MLP.
GFLOP/s $=$ FLOPs totales$/T_p$.

\section{Trampas que caza el profe}

\begin{itemize}
  \item Decir \textbf{``broadcast''} en memoria compartida $\to$ di
        \textbf{``lectura concurrente CREW''}.
  \item Confundir tiempo paralelo con \textbf{suma} de tiempos $\to$ es el
        \textbf{máximo} (worker crítico).
  \item Decir que $E(p)$ siempre es 1 $\to$ cae con $p$ (solo se mantiene en
        weak scaling).
  \item Decir ``usamos OpenMP'' sin aclarar $\to$ ``modelo OpenMP, implementado
        con multiprocessing por el GIL''.
\end{itemize}

% ====================================================================
\part*{PARTE X --- Chuleta final}
\addcontentsline{toc}{section}{PARTE X --- Chuleta final}

\section{Diccionario rápido}

\begin{tabular}{ll}
\toprule
Término & Significado \\\midrule
$T_s$ & tiempo secuencial (con 1 procesador) \\
$T_p$ & tiempo paralelo (con $p$ procesadores) \\
$S(p)=T_s/T_p$ & speedup (cuántas veces más rápido) \\
$E(p)=S(p)/p$ & eficiencia (aprovechamiento, $\le 1$) \\
$W$ & work = trabajo total = $O(T_s)$ \\
$T_\infty$ & span = camino crítico (tiempo con $\infty$ procesadores) \\
$C_p=T_p\cdot p$ & costo \\
PRAM & modelo teórico de máquina paralela \\
CREW & lectura Concurrente, Escritura Exclusiva \\
\texttt{pardo} & bucle en paralelo \\
Brent & $T_p\le W/p + T_\infty$ \\
Amdahl & techo de speedup por fracción secuencial \\
Gustafson & speedup casi lineal (weak scaling) \\
GIL & candado de Python que impide hilos paralelos \\
fork & crear proceso hijo que hereda memoria \\
MLP & perceptrón multicapa (red neuronal) \\
$\eta$ & learning rate \\
$E$ & épocas \\
accuracy & fracción de aciertos \\
$p_s$ & número de semillas (32) \\
\bottomrule
\end{tabular}

\section{Las fórmulas que debes saber}

\[
S(p)=\frac{T_s}{T_p}\qquad E(p)=\frac{S(p)}{p}\qquad
T_p\le\frac{W}{p}+T_\infty\qquad T_p\approx a\frac{p_s}{p}n+b\log_2 p
\]

\[
\text{FLOPs por MLP}=4Eh(d+c)n\qquad
\text{GFLOP/s}=\frac{\text{FLOPs totales}}{T_p\cdot 10^9}
\]

\section{Los números del proyecto (memorízalos)}

\begin{itemize}
  \item $p_s=32$ semillas, $d=32$, $h=128$, $c=10$, $E=100$ épocas.
  \item Dos plataformas: MacBook M5 (10 núcleos) y \textbf{cluster Khipu n003}
        (Xeon Gold 6130, 32 núcleos lógicos $\approx$ 16 físicos con HT).
  \item Speedup techo: \textbf{$\sim$5 (Mac), $\sim$10.6 (Khipu)} --- crece con
        el hardware.
  \item Reducción arg-máx $<1.5\%$ del tiempo (embarassingly parallel).
  \item Ajuste teórico (Khipu): $a=2.42\!\times\!10^{-4}$\,s/muestra,
        $b=2.18$\,s, $R^2=0.989$.
  \item GFLOP/s pico: $\sim$53 (Mac), \textbf{$\sim$76 (Khipu)}.
  \item Por qué no llega a $p$: hyper-threading (32 lógicos $\neq$ 32 físicos) +
        contención de memoria (MLP memory-bound).
\end{itemize}

\begin{caja}[title=\textbf{Regla de oro}]
Si el profe pregunta algo que no medimos, di ``no lo medimos, pero\ldots'' y da
la respuesta cualitativa correcta. \textbf{Nunca inventes números.}
\end{caja}

\end{document}
